\begin{abstract}
多智能体语言模型系统通常优化最终聚合结果，却很少观测产生该结果的行为。当完整独立报告比短答案或模仿具有更高成本时，系统无法区分真实审议与低成本一致。本文提出\emph{AEIC-MD}：以假设--证据--推理--结论（AEIC）报告作为机制设计者与LLM Agent之间可观测接口的审议框架。求解前，每个Agent观察公开信用--成本规则，并从低资源报告、完整AEIC推理、追加验证和同伴模仿中选择行动；行动控制真实生成过程，并在结算前写入日志，从而可以直接检验追随者响应。

在StrategyQA、MMLU-Hard和TruthfulQA的450道配对测试题上，静态AEIC规则将高资源行动率从0.46\%提高到56.93\%，多数投票准确率从68.00\%提高到82.22\%（$+14.22$个百分点；Holm校正McNemar $p=1.23\times10^{-11}$）。仅结果奖励只将高资源行动率提高到3.08\%，且未改善准确率（67.33\%，$p=0.780$）。奖励和成本扫描呈现预期单调响应；受控信息环境、额外Seed、信号与历史消融以及3--5 Agent面板限定了稳健性。留出GLM后端未复现行动或准确率增益（69.33\%对70.67\%，$p=0.791$），表明合同必须针对后端校准。八组预注册实验共包含5,850条经核验的问题级记录。结果支持有边界的结论：针对可验证报告结构的公开规则能够改变追随者资源行动并改善集体决策，但合同语义不会自动跨模型家族迁移。
\keywords{多智能体大语言模型 \and 机制设计 \and 结构化推理 \and 群体审议 \and 追随者响应}
\end{abstract}
